\section{Related Work}

\noindent\textbf{Quantum Simulation and NISQ Systems.}
Computer architecture uses measurement and modeling to study immature systems~\cite{skadron2003challenges}; \SystemName applies that discipline to quantum systems. ScaleQsim, AURORA-Q, and \cuQuantum improve state-vector simulation through distribution, GPUs, tiered memory, asynchronous movement, and optimized primitives~\cite{scaleqsim,auroraq,cuquantum}; NISQ studies show current-device limits~\cite{preskill2018quantum}. SupermarQ, QASMBench, MQT Bench, QAOAKit, QED-C, and QCHALLenge provide circuits, metrics, parameters, and application-aware tests~\cite{supermarq,qasmbench,mqtbench,qaoakit,qedcbench,qchallenge}. \SystemName asks a different question: what hardware and quality threshold lets a quantum-circuit application beat a native HPC application?

\noindent\textbf{Quantum Applications and Native Baselines.}
Quantum kernels, variational classifiers, VQE chemistry, QAOA optimization, and Hamiltonian simulation are common practical-advantage candidates~\cite{biamonte2017quantum,cerezo2022challenges,kandala2017hardware,farhi2014qaoa,georgescu2014quantum,daley2022practical}, but each domain also has strong native baselines and different quality metrics~\cite{scikit-learn}. Fault-tolerant and resource-estimation work shows that logical cost depends on code distance, cycle time, layout, magic-state factories, decoder latency, and physical assumptions~\cite{litinski2019game,azureestimator,webber2022impact}; hybrid runtimes add queue, fidelity, and classical-resource constraints~\cite{qurator,k8squantum}. \SystemName supplies the same-task baselines, quality targets, and threshold artifacts these studies need.
